\documentclass[a4paper,11pt]{article}
\usepackage[T1]{fontenc}
\usepackage[utf8]{inputenc}
\usepackage[ngerman]{babel}
\usepackage[margin=2.5cm]{geometry}

\begin{document}

\section*{Blatt 3, Aufgabe 3.1.4: Histogramme}

\textbf{Gleichverteilung.}
Mit größerem Variationskoeffizienten (CV) verteilen sich die Werte über ein
breiteres Intervall. Die Balken im Inneren sind ungefähr gleich hoch. Der
auffällig hohe erste Balken bei CV~1 und CV~2 entsteht durch die negativen
Werte. Sie machen \(21{,}04\,\%\) bzw. \(35{,}53\,\%\) der Stichprobe aus und
werden dem ersten Interval hinzugerechnet.

\medskip
\noindent\textbf{Exponential- und Erlang-Verteilung.}
Bei CV~1 fallen die Balken bei beiden Verteilungen von links nach rechts ab.
Das passt dazu, dass die Erlang-Verteilung mit \(k=1\) eine
Exponentialverteilung ist. Bei \(k=100\) liegen die Werte dagegen dicht um den
Mittelwert 1. Das Histogramm sieht dort annähernd glockenförmig aus. Dies deckt sich mit der möglichen Aproximation durch die Normalverteilung.

\medskip
\noindent\textbf{Hyperexponentialverteilung \(H_2\).}
Bei CV~1 sind beide Raten gleich, daher hat das Histogramm wieder die Form
einer Exponentialverteilung. Bei CV~2 sieht man mehr sehr kleine Werte und
zugleich deutlich mehr große Werte. Rund \(1{,}19\,\%\) der Ziehungen liegen
bei 10 oder darüber. Bei der Exponentialverteilung sind es nur
\(0{,}0042\,\%\). Die Werte ab 10 landen im letzen Intervall, dementsprechend der Peak.


\end{document}
